% ============================================================
%  Dok 257 | Scale-Dependent Information Units as
%           Phase Transitions on T⁴ — Version 2
%  Chapter content EN
% ============================================================

\title{Scale-Dependent Information Units as\\
Phase Transitions on \(T^4\)\\[0.4em]
\large Geometry as the Primary Foundation of Information and Energy}
\author{Johann Pascher}
\date{May 2026}
\maketitle

\begin{abstract}
Dok.~255 derived Landauer's principle and Vopson's mass-information
equivalence from \(T^4\) geometry, identifying information with
topologically frozen kinetic energy.
The present refinement shows that this derivation holds as a
special case at the thermodynamic scale, not universally.
The more fundamental picture is: the minimal information unit is
by definition the topological winding quantum \(\Delta w = 1\) —
scale-universal and energy-independent.
The corresponding energy \(E_\text{bit}(L) = \hbar c / L\)
follows from action and the uncertainty relation and is scale-specific.
Landauer's \(E_L = k_B T \ln 2\) is the special case at the
temperature-length scale where \(E_\text{bit}(L) = k_BT\ln2\) —
biological room temperature.
All phase transitions are discrete because
\(\pi_1(T^4) \cong \mathbb{Z}^4\) admits no half-integer winding quanta.
Geometry is more fundamental than energy.
We work in natural units with \(c = \hbar = k_B = 1\)
unless stated otherwise.
\end{abstract}

\tableofcontents
\newpage

% ------------------------------------------------------------
\section{The Problem: Three Levels, One Hierarchy}
% ------------------------------------------------------------

Dok.~255 stated:
\begin{center}
\textit{Information is topologically frozen kinetic energy.}
\end{center}

This is correct as a description of the thermodynamic regime.
The deeper question is: what is more fundamental — geometry or energy?

The internal consistency check shows that the equivalence temperature
at which \(E_\text{bit}(L) = k_B T \ln 2\) holds
varies over \textbf{70 orders of magnitude}:

\begin{center}
\renewcommand{\arraystretch}{1.4}
\begin{tabular}{lrr}
\toprule
\textbf{Scale} & \textbf{L [m]} & \textbf{\(T_\text{equiv}\) [K]} \\
\midrule
Planck         & \(1.6 \times 10^{-35}\) & \(2.0 \times 10^{32}\) \\
FFGFT \(L_0\)  & \(2.1 \times 10^{-39}\) & \(1.5 \times 10^{36}\) \\
QCD/Proton     & \(1.0 \times 10^{-15}\) & \(3.3 \times 10^{12}\) \\
Atom           & \(1.0 \times 10^{-10}\) & \(3.3 \times 10^{7}\)  \\
DNA codon      & \(1.0 \times 10^{-9}\)  & \(3.3 \times 10^{6}\)  \\
\textbf{Cell}  & \(\mathbf{1.0 \times 10^{-5}}\) & \(\mathbf{331}\) \\
Cosmic         & \(4.4 \times 10^{26}\) & \(7.5 \times 10^{-30}\) \\
\bottomrule
\end{tabular}
\end{center}

with the formula (cf.\ Section~\ref{sec:three}):
\begin{equation}
  T_\text{equiv}(L) = \frac{\hbar c}{k_B \, L \, \ln 2}.
  \label{eq:Tequiv}
\end{equation}

Only at the cellular scale does \(T_\text{equiv} \approx 331\) K
match biological room temperature.

% ------------------------------------------------------------
\section{Three Independent Levels}
\label{sec:three}
% ------------------------------------------------------------

The calculation enforces a separation into three levels:

\paragraph{Level 1: Topological (universal).}

\begin{definition}[Minimal information unit]
The minimal information unit on \(T^4\) is the winding quantum
\(I_\text{bit} = \Delta w = 1\), \(\Delta w \in \mathbb{Z}\).
One bit is identified with a minimal topological winding difference
\(\Delta w = 1\).
\end{definition}

This identification is the conceptual bridge between
topology and information theory.
The unit is scale-independent: \(\pi_1(T^4) \cong \mathbb{Z}^4\)
enforces discrete integer winding numbers at every scale.
No phase transition can produce half a winding quantum.

\paragraph{Level 2: Geometric (scale-specific).}

The energy of a winding quantum at length scale \(L\) follows
from action and the uncertainty relation:
The action of one winding quantum is \(S = \Delta w \cdot h = h\).
From the uncertainty relation \(\Delta x \cdot \Delta p \sim \hbar\)
at spatial scale \(L\) one obtains \(p \sim \hbar/L\), hence:
\begin{equation}
  E_\text{bit}(L) = pc = \frac{\hbar c}{L}.
  \label{eq:Ebit}
\end{equation}

\textit{Note:} The numerical prefactor is a scale definition;
the formula gives the correct order of magnitude.
In natural units (\(c = \hbar = k_B = 1\)), \(E_\text{bit}(L) = 1/L\).

\paragraph{Level 3: Thermodynamic (one scale).}

Landauer's principle \cite{Landauer1961}:
\begin{equation}
  E_L = k_B T \ln 2 \quad \text{per bit}
  \label{eq:Landauer}
\end{equation}
holds in the thermodynamic regime, i.e.\ when
\(E_\text{bit}(L) = E_L\).
Setting equal yields equation~(\ref{eq:Tequiv}).

\subsection{Visualisation of the Hierarchy}

\begin{center}
\renewcommand{\arraystretch}{1.6}
\begin{tabular}{clll}
\toprule
\textbf{Level} & \textbf{Character} & \textbf{Quantity} & \textbf{Validity} \\
\midrule
1 & topological   & \(\Delta w = 1\)                & universal, all scales \\
  & \(\downarrow\) & & \\
2 & geometric     & \(E_\text{bit}(L) = \hbar c/L\) & every scale \(L\) \\
  & \(\downarrow\) & & \\
3 & thermodynamic & \(E_L = k_BT\ln2\)             & only when \(E_\text{bit} = k_BT\) \\
\bottomrule
\end{tabular}
\end{center}

% ------------------------------------------------------------
\section{Geometry is More Fundamental than Energy}
% ------------------------------------------------------------

The derivation hierarchy is:
\begin{align}
  T^4\text{-geometry} \;&\longrightarrow\;
  \Delta w = 1 \;\longrightarrow\;
  E_\text{bit}(L) = \hbar c / L \notag\\
  &\longrightarrow\;
  E_L = k_BT\ln 2 \text{ (special case)}.
\end{align}

Energy is a scale-dependent projection of \(T^4\) geometry,
not its foundation.

\begin{center}
\renewcommand{\arraystretch}{1.4}
\begin{tabular}{lp{8cm}}
\toprule
\textbf{Dok.~255} & Information = frozen kinetic energy \\
\textbf{Dok.~257} & Information = winding quantum \(\Delta w\);
                    energy is a scale-specific projection of it \\
\midrule
\multicolumn{2}{l}{\textit{Both statements consistent:
Dok.~255 describes the thermodynamic special case.}} \\
\bottomrule
\end{tabular}
\end{center}

% ------------------------------------------------------------
\section{Phase Transitions are Discrete at Every Scale}
% ------------------------------------------------------------

All phase transitions in physics are discrete.
This follows directly from \(\pi_1(T^4) \cong \mathbb{Z}^4\).

\begin{center}
\renewcommand{\arraystretch}{1.4}
\begin{tabular}{lll}
\toprule
\textbf{Scale} & \textbf{Phase transition} & \textbf{Information unit} \\
\midrule
Planck/FFGFT & Topological enclosure     & \(\Delta w=1\), \(E \sim E_P\) \\
QCD          & Confinement, hadronisation & Baryon \(B=1\), \(E\sim\Lambda_\text{QCD}\) \\
Atom         & Electron configuration    & Quantum number \(n\), \(E\sim\) eV \\
Molecule/DNA & Codon, folding            & 3 base pairs, \(E\sim k_BT\) \\
Cell         & Protein folding           & Domain, \(E\sim k_BT\) \\
Cosmic       & \(\sigma\)-shift          & \(\hbar c/R_U\), \(E \ll k_BT_\text{CMB}\) \\
\bottomrule
\end{tabular}
\end{center}

At every scale: \(I_\text{bit} = \Delta w = 1\) (topological).
The energy of the bit varies with scale.

% ------------------------------------------------------------
\section{Scale Resonance at the Scale of Life — Hypothesis}
% ------------------------------------------------------------

For \(L \sim 10^{-5}\) m (cell size) and \(T \sim 300\) K:
\begin{equation}
  E_\text{bit}(L_\text{cell}) = \frac{\hbar c}{10^{-5}\,\text{m}}
  \approx 3.2 \times 10^{-21}\,\text{J}
  \approx k_B \cdot 331\,\text{K} \cdot \ln 2.
\end{equation}

This suggests the following hypothesis:

\begin{center}
\textit{Life as a self-organising, information-processing structure\\
may only exist at scales where \(E_\text{bit}(L) \approx k_BT\).\\
Biological room temperature would then be a consequence of this scale resonance.}
\end{center}

Whether this coincidence is causal or accidental remains open
and deserves further investigation.
It is stated here as a hypothesis, not a derivation.

% ------------------------------------------------------------
\section{Refinement of Dok.~255}
% ------------------------------------------------------------

\begin{enumerate}
  \item \textbf{Correct in Dok.~255:}
        Landauer and Vopson follow from \(T^4\) geometry
        in the thermodynamic regime.
  \item \textbf{Refinement:}
        The universal information unit is \(\Delta w = 1\)
        (topological), not \(k_BT\ln2\) (thermodynamic).
  \item \textbf{Extension:}
        Energy is a scale-specific projection of geometry.
  \item \textbf{Consequence:}
        \(I_\text{bit} = \Delta w\) holds at all scales.
        \(E_L = k_BT\ln2\) holds only at the scale of life.
        \(E_\text{bit} = \hbar c/L\) holds geometrically at every scale.
\end{enumerate}

% ------------------------------------------------------------
\section{Summary}
% ------------------------------------------------------------

\begin{enumerate}
  \item The minimal information unit is by definition the winding
        quantum \(\Delta w = 1\) — universal from
        \(\pi_1(T^4) \cong \mathbb{Z}^4\).
  \item The energy of a bit follows from \(S = h\) and
        \(\Delta x\Delta p \sim \hbar\):
        \(E_\text{bit}(L) = \hbar c / L\).
  \item Landauer \(E_L = k_BT\ln2\) is the special case
        at biological room temperature (\(L \sim 10^{-5}\) m,
        \(T \sim 300\) K).
  \item All phase transitions are discrete because
        \(\pi_1(T^4) \cong \mathbb{Z}^4\) is integer-valued.
  \item Geometry is more fundamental than energy:
        \(T^4 \to \Delta w \to E_\text{bit}(L) \to E_L\) (special case).
  \item Scale resonance at the scale of life is an open hypothesis,
        not a derivation.
\end{enumerate}

\bigskip
\noindent
\textbf{Methodological classification:}
Dok.~257 is an algebraic bridge from \(T^4\) and
\(\pi_1(T^4) \cong \mathbb{Z}^4\).
The scale resonance hypothesis requires further investigation.

\bigskip
\noindent
\textbf{Integration in the corpus.}
The three-level hierarchy formulated here is the field-theoretically
developed form of the static/dynamic demarcation of Doc.~261:
level~1 (topological, \(\Delta w\)) is layer~1 -- the static,
reference-free \(\xi\)-geometry level. Levels~2 and~3
(\(E_\text{bit}(L) = \hbar c/L\); \(E_L = k_B T \ln 2\)) are layer-2
projections and therefore carry the embedding constants
\(\hbar, c, k_B\). The geometric formula \(E_\text{bit}(L) = \hbar c/L\)
coincides with the general energy hierarchy \(E_N = \hbar c/L_N\) of
Doc.~184. The Vopson factor
\(\xi_\text{FFGFT}/\xi_\text{UIFT} \approx 414{,}684\) discussed in
Doc.~013 appears, in the present picture, as a projection ratio between
a layer-1 quantity (\(\xi_\text{FFGFT}\), dimensionless) and a layer-2
quantity (\(\xi_\text{UIFT} = k_B T_\text{CMB}\ln 2/(m_e c^2)\), bound
to the CMB temperature).

\bigskip
\begin{flushright}
\textit{Johann Pascher, Upper Austria, May 2026}
\end{flushright}

\begin{thebibliography}{99}

\bibitem{Pascher255}
J.~Pascher,
\textit{Information as Topologically Frozen Kinetic Energy},
Dok.~255, Zenodo DOI: 10.5281/zenodo.20117635 (2026).

\bibitem{Landauer1961}
R.~Landauer,
\textit{Irreversibility and heat generation in the computing process},
IBM Journal of Research and Development \textbf{5}(3), 183--191 (1961).

\bibitem{Berut2012}
A.~Bérut et al.,
\textit{Experimental verification of Landauer's principle},
Nature \textbf{483}, 187--189 (2012).

\bibitem{Vopson2019}
M.\,M.~Vopson,
\textit{The mass-energy-information equivalence principle},
AIP Advances \textbf{9}(9), 095206 (2019).

\end{thebibliography}
